\ExerciseSection

\begin{enumerate}
\def\labelenumi{\arabic{enumi})}
\tightlist
\item
  设 \(\varphi\) 是 \(\mathbf{R}^n\) 到 \(\mathbf{T}^n\) 上的典范同态，\(u\) 是 \(\mathbf{R}^m\) 到 \(\mathbf{R}^n\) 中的一个线性映射。要使 \(\varphi \circ u\) 是 \(\mathbf{R}^m\) 到 \(\mathbf{T}^n\) 中的一个严格态射，必须且只须下列条件之一成立：
\end{enumerate}

\begin{enumerate}
\def\labelenumi{\alph{enumi})}
\item
  \(\overline{u}^1 (\mathbf{Z}^n)\) 的秩为 \(m\) ；或者
\item
  \(u(\mathbf{R}^m)\cap \mathbf{Z}^n\) 的秩等于 \(u(\mathbf{R}^m)\) 的维数。
\end{enumerate}

\begin{enumerate}
\def\labelenumi{\arabic{enumi})}
\setcounter{enumi}{1}
\item
  证明：从 p 进数的加法拓扑群 \(Q_{p}\) （第三章，§ 6，习题 23 及其后各题）到加法群 R 中的连续同态只有零映射。
\item
  每个 \(p\) 进数 \(x\) （习题 2）都可以写成形式 \(\sum_{-\infty}^{+\infty} \alpha_k p^k\) ，其中诸 \(\alpha_k\) 是有理整数，且指标 \(k < 0\) 的那些除对指标 \(k\) 的有限多个值外全为零。此外，若 \(\sum_{-\infty}^{+\infty} \alpha_k p^k = \sum_{-\infty}^{+\infty} \beta_k p^k\) ，则有理数 \(\sum_{-\infty}^{-1} \alpha_k p^k\) 与 \(\sum_{-\infty}^{-1} \beta_k p^k\) 模 1 同余。设 \(\varphi_p(x)\) 是有理数 \(\sum_{-\infty}^{-1} \alpha_k p^k\) 的模 1 类。\\
\end{enumerate}

\begin{enumerate}
\def\labelenumi{\alph{enumi})}
\tightlist
\item
  证明 \(\varphi_p\) 是拓扑群 \(\mathbf{Q}_p\) 到 \(\mathbf{T}\) 中的一个连续同态，并且对每个连续同态 \(f\) （\(\mathbf{Q}_p\) 到 \(\mathbf{T}\) 中的），存在唯一的 \(a \in \mathbf{Q}_p\) ，使得 \(f(x) = \varphi_p(ax)\) 对一切 \(x \in \mathbf{Q}_p\) 成立。{[}注意，知道了诸元素 \(f(p^k)\) （\(k \leqslant 0\) ）就唯一地确定了 \(f\) ；证明这些元素中的每一个都是分母为 \(p\) 的幂的分数的模 1 类；由此证明存在 \(a \in \mathbf{Q}_p\) ，使得 \(f(p^k) = \varphi_p(ap^k)\) 对一切 \(k \leqslant 0\) 成立。{]}\\
\item
  证明：若 \(\varphi: \mathbf{R} \to \mathbf{T}\) 是典范同态，则 \(\varphi(x) = \sum_{p} \varphi_p(x)\) 对一切 \(x \in \mathbf{Q}\) 成立，其中和式取遍所有素数 \(p\) （注意该和式中除有限多项外全为零）。
\end{enumerate}

¶ 4) a) 设 G 是一个局部紧交换群，H 是 G 的一个同构于 \(\mathbf{R}^{p} \times \mathbf{T}^{q}\) 的子群，其中 \(p, q\) 是 \(\geqslant 0\) 的整数；再设 G/H 同构于 R 或 T 之一。证明 G 同构于 H 与一个同构于 G/H 的闭子群 L 的乘积。{[}应用第五章 § 3 习题 7 d)，得到 R 到 G 中的一个连续同态 \(f\) ，使得 \(f(\mathbf{R}) \subset \mathbf{H}\) ；由此构造另一个连续同态 \(g: \mathbf{R} \to \mathbf{G}\) ，使得 \(f(t) - g(t) \in \mathbf{H}\) 对一切 \(t \in \mathbf{R}\) 成立，并且 \(g(\mathbf{R}) \cap \mathbf{H} = \{e\}\) 。注意 \(f(\mathbf{H})\) 是 R 的一个真闭子群，并且对 H 中每个元素 \(z \neq e\) ，存在 R 到 H 中的一个连续同态 \(u\) ，使得 \(u(1) = z\) 。{]}

\begin{enumerate}
\def\labelenumi{\alph{enumi})}
\setcounter{enumi}{1}
\tightlist
\item
  设 G 是一个豪斯多夫交换拓扑群，H 是 G 的一个同构于 \(R^{p} \times T^{q}\) 的闭子群，并设存在一个满射连续同态 \(\varphi: R \to G/H\) 。证明存在连续同态 \(f: R \to G\) ，使得 \(f(\mathbf{R}) \cap \mathbf{H} = \{e\}\) 且 \(\mathbf{G} = \mathbf{H}. f(\mathbf{R})\) 。{[}通过考虑给定拓扑与这样一个拓扑在 G 上的上确界——按该拓扑，o 的一个基本
\end{enumerate}

邻域系由诸 \(\overline{p}(\varphi(I))\) 组成，其中 \(p:G\to G/H\) 是典范映射，I 取遍 R 中 o 的一个基本邻域系——化归为情形 a)；证明 G 在这个新拓扑中是局部紧的（参看第三章，§ 4，习题 10 e）。{]}

\begin{enumerate}
\def\labelenumi{\alph{enumi})}
\setcounter{enumi}{2}
\tightlist
\item
  给出一个例子，使得 \(p\) 在 \(f(\mathbf{R})\) 上的限制不是双连续的{[}参看 § 1，习题 2f）{]}。
\end{enumerate}

\begin{enumerate}
\def\labelenumi{\arabic{enumi})}
\setcounter{enumi}{4}
\item
  设 G 是一个连通拓扑群，H 是 G 的一个没有任意小的子群的紧正规子群（第三章，§ 2，习题 30）。证明 H 包含在 G 的中心内。（注意，对 G 中充分接近 e 的 s，H 在 \(x \rightarrow sxs^{-1}\) 下的像任意接近于 e。）
\item
  设 \(\mathbf{G}\) 是一个拓扑群，\(\mathbf{H}\) 是 \(\mathbf{G}\) 的一个闭正规子群，它同构于 \(\mathbf{R}^n (n\geqslant 1)\) ，并且使得 \(\mathbf{G} / \mathbf{H}\) 是交换的。
\end{enumerate}

\begin{enumerate}
\def\labelenumi{\alph{enumi})}
\item
  再设 H 除 H 自身与 \(\{e\}\) 外不含 G 的连通正规子群，并且 H 不包含在 G 的中心内。给定一个不与 H 的每个元素都交换的 \(x_{0} \in G\) ，证明：对每个 \(v \in H\) ，存在唯一的 \(u \in H\) 使得 \(x_{0}^{-1}ux_{0}u^{-1} = v\) 。（注意，在 H 的加法记号下，该方程变为 \(u - x_{0}^{-1}ux_{0} = -v\) ，并且 H 的连续自同态 \(u \to u - x_{0}^{-1}ux_{0}\) 是 \(R^{n}\) 到其自身的一个线性映射。证明它的核是 G 的一个正规子群。）此外，若 \(f(u)\) 是唯一的 \(u \in H\) ，它使 \(x_{0}^{-1}ux_{0}u^{-1} = v\) ，则 f 是 H 的一个双连续自同构。
\item
  在 a) 的假设下，证明 G 到其自身的连续映射 \(g: y \to (f(x_{0}^{-1}yx_{0}y^{-1}))^{-1}y\) 以 G 中由所有与 \(x_{0}\) 交换的元素组成的子群 L 为像，并且关系 \(g(y) = g(y')\) 等价于 \(yy'^{-1} \in H\) 。由此推出存在一个连续双射 \(h: G/H \to L\) ，它的逆是典范映射 \(p: G \to G/H\) 在 L 上的限制，进而推出 G 同构于 H 与 L 的拓扑半直积。
\end{enumerate}

\begin{enumerate}
\def\labelenumi{\arabic{enumi})}
\setcounter{enumi}{6}
\tightlist
\item
  设 G 是一个拓扑群，H 是一个同构于 \(R^{p} \times T^{q}\) 的正规子群（ \(p \geqslant 0\) ，\(q \geqslant 0\) ），并设存在一个满射连续同态 \(\varphi: R \to G/H\) 。证明存在一个连续同态 \(f: R \to G\) ，使得 \(f(R) \cap H = \{e\}\) 且 \(G = H\) . \(f(R)\) 。{[}利用习题 4 b) 与第五章 § 3 习题 8，先化归为 G 不是交换群的情形，再利用习题 5 化归为 q = 0 的情形；然后对 p 用归纳法，并结合习题 6。{]}
\end{enumerate}

